% Options for packages loaded elsewhere
\PassOptionsToPackage{unicode}{hyperref}
\PassOptionsToPackage{hyphens}{url}
\documentclass[
  12pt,
]{article}
\usepackage{xcolor}
\usepackage{amsmath,amssymb}
\setcounter{secnumdepth}{-\maxdimen} % remove section numbering
\usepackage{iftex}
\ifPDFTeX
  \usepackage[T1]{fontenc}
  \usepackage[utf8]{inputenc}
  \usepackage{textcomp} % provide euro and other symbols
\else % if luatex or xetex
  \usepackage{unicode-math} % this also loads fontspec
  \defaultfontfeatures{Scale=MatchLowercase}
  \defaultfontfeatures[\rmfamily]{Ligatures=TeX,Scale=1}
\fi
\usepackage{lmodern}
\ifPDFTeX\else
  % xetex/luatex font selection
\fi
% Use upquote if available, for straight quotes in verbatim environments
\IfFileExists{upquote.sty}{\usepackage{upquote}}{}
\IfFileExists{microtype.sty}{% use microtype if available
  \usepackage[]{microtype}
  \UseMicrotypeSet[protrusion]{basicmath} % disable protrusion for tt fonts
}{}
\makeatletter
\@ifundefined{KOMAClassName}{% if non-KOMA class
  \IfFileExists{parskip.sty}{%
    \usepackage{parskip}
  }{% else
    \setlength{\parindent}{0pt}
    \setlength{\parskip}{6pt plus 2pt minus 1pt}}
}{% if KOMA class
  \KOMAoptions{parskip=half}}
\makeatother
\usepackage{color}
\usepackage{fancyvrb}
\newcommand{\VerbBar}{|}
\newcommand{\VERB}{\Verb[commandchars=\\\{\}]}
\DefineVerbatimEnvironment{Highlighting}{Verbatim}{commandchars=\\\{\}}
% Add ',fontsize=\small' for more characters per line
\newenvironment{Shaded}{}{}
\newcommand{\AlertTok}[1]{\textcolor[rgb]{1.00,0.00,0.00}{\textbf{#1}}}
\newcommand{\AnnotationTok}[1]{\textcolor[rgb]{0.38,0.63,0.69}{\textbf{\textit{#1}}}}
\newcommand{\AttributeTok}[1]{\textcolor[rgb]{0.49,0.56,0.16}{#1}}
\newcommand{\BaseNTok}[1]{\textcolor[rgb]{0.25,0.63,0.44}{#1}}
\newcommand{\BuiltInTok}[1]{\textcolor[rgb]{0.00,0.50,0.00}{#1}}
\newcommand{\CharTok}[1]{\textcolor[rgb]{0.25,0.44,0.63}{#1}}
\newcommand{\CommentTok}[1]{\textcolor[rgb]{0.38,0.63,0.69}{\textit{#1}}}
\newcommand{\CommentVarTok}[1]{\textcolor[rgb]{0.38,0.63,0.69}{\textbf{\textit{#1}}}}
\newcommand{\ConstantTok}[1]{\textcolor[rgb]{0.53,0.00,0.00}{#1}}
\newcommand{\ControlFlowTok}[1]{\textcolor[rgb]{0.00,0.44,0.13}{\textbf{#1}}}
\newcommand{\DataTypeTok}[1]{\textcolor[rgb]{0.56,0.13,0.00}{#1}}
\newcommand{\DecValTok}[1]{\textcolor[rgb]{0.25,0.63,0.44}{#1}}
\newcommand{\DocumentationTok}[1]{\textcolor[rgb]{0.73,0.13,0.13}{\textit{#1}}}
\newcommand{\ErrorTok}[1]{\textcolor[rgb]{1.00,0.00,0.00}{\textbf{#1}}}
\newcommand{\ExtensionTok}[1]{#1}
\newcommand{\FloatTok}[1]{\textcolor[rgb]{0.25,0.63,0.44}{#1}}
\newcommand{\FunctionTok}[1]{\textcolor[rgb]{0.02,0.16,0.49}{#1}}
\newcommand{\ImportTok}[1]{\textcolor[rgb]{0.00,0.50,0.00}{\textbf{#1}}}
\newcommand{\InformationTok}[1]{\textcolor[rgb]{0.38,0.63,0.69}{\textbf{\textit{#1}}}}
\newcommand{\KeywordTok}[1]{\textcolor[rgb]{0.00,0.44,0.13}{\textbf{#1}}}
\newcommand{\NormalTok}[1]{#1}
\newcommand{\OperatorTok}[1]{\textcolor[rgb]{0.40,0.40,0.40}{#1}}
\newcommand{\OtherTok}[1]{\textcolor[rgb]{0.00,0.44,0.13}{#1}}
\newcommand{\PreprocessorTok}[1]{\textcolor[rgb]{0.74,0.48,0.00}{#1}}
\newcommand{\RegionMarkerTok}[1]{#1}
\newcommand{\SpecialCharTok}[1]{\textcolor[rgb]{0.25,0.44,0.63}{#1}}
\newcommand{\SpecialStringTok}[1]{\textcolor[rgb]{0.73,0.40,0.53}{#1}}
\newcommand{\StringTok}[1]{\textcolor[rgb]{0.25,0.44,0.63}{#1}}
\newcommand{\VariableTok}[1]{\textcolor[rgb]{0.10,0.09,0.49}{#1}}
\newcommand{\VerbatimStringTok}[1]{\textcolor[rgb]{0.25,0.44,0.63}{#1}}
\newcommand{\WarningTok}[1]{\textcolor[rgb]{0.38,0.63,0.69}{\textbf{\textit{#1}}}}
\usepackage{longtable,booktabs,array}
\usepackage{calc} % for calculating minipage widths
% Correct order of tables after \paragraph or \subparagraph
\usepackage{etoolbox}
\makeatletter
\patchcmd\longtable{\par}{\if@noskipsec\mbox{}\fi\par}{}{}
\makeatother
% Allow footnotes in longtable head/foot
\IfFileExists{footnotehyper.sty}{\usepackage{footnotehyper}}{\usepackage{footnote}}
\makesavenoteenv{longtable}
\setlength{\emergencystretch}{3em} % prevent overfull lines
\providecommand{\tightlist}{%
  \setlength{\itemsep}{0pt}\setlength{\parskip}{0pt}}
\usepackage{bookmark}
\IfFileExists{xurl.sty}{\usepackage{xurl}}{} % add URL line breaks if available
\urlstyle{same}
\hypersetup{
  hidelinks,
  pdfcreator={LaTeX via pandoc}}

\author{SCX}
\date{2026}

\begin{document}

\section{Theorem 2: Weak Feature Failure Lower
Bound}\label{theorem-2-weak-feature-failure-lower-bound}

\begin{quote}
when feature表示 \(\phi(x)\)
包含'strue state信息non-足时，基于一致性'sNoiseDetectionMethod（包括
SCX）No法优于损失基线。本定理FromInformation Theory角度量化了this 一边界conditions。
\end{quote}

\textbf{关联Experiment}: DermaMNIST (SCX-Health), AlN v3 MLIP (12-dim
手工feature) \textbf{Related Code}: \texttt{src/scx/state/discovery.py},
\texttt{src/scx/valuation/noise\_score.py},
\texttt{src/scx/expert/reliability.py} \textbf{Related Concepts}:
{[}{[}弱feature失效{]}{]}, {[}{[}State-Conditioned eXpertise{]}{]}

\begin{center}\rule{0.5\linewidth}{0.5pt}\end{center}

\subsection{Table of Contents}\label{ux76eeux5f55}

\begin{enumerate}
\def\labelenumi{\arabic{enumi}.}
\tightlist
\item
  \hyperref[1-ux7b26ux53f7ux4e0eux8bbeux5b9a]{Symboland 设定}
\item
  \hyperref[2-ux5b9aux4e49ux5f31ux7279ux5f81]{定义：弱feature}
\item
  \hyperref[3-ux5f15ux7406-1ux72b6ux6001ux4f30ux8ba1ux8befux5deeux4e0bux754c]{引理
  1：state estimation误差lower bound}
\item
  \hyperref[4-ux5f15ux7406-2scx-ux4f30ux8ba1ux9000ux5316]{引理 2：SCX
  估计退化}
\item
  \hyperref[5-ux5b9aux7406-2ux5f31ux7279ux5f81ux5931ux8d25ux4e0bux754c]{定理
  2：Weak Feature Failurelower bound}
\item
  \hyperref[6-ux63a8ux8bba]{Corollaries}
\item
  \hyperref[7-ux5b9eux7528ux89e3ux8bfb]{实using 解读}
\item
  \hyperref[8-ux4e0e-theorem-1-ux7684ux8054ux7cfb]{and  Theorem 1 's联系}
\end{enumerate}

\begin{center}\rule{0.5\linewidth}{0.5pt}\end{center}

\subsection{1. Symboland 设定}\label{ux7b26ux53f7ux4e0eux8bbeux5b9a}

\subsubsection{1.1 Data Generation}\label{ux6570ux636eux751fux6210}

\begin{longtable}[]{@{}
  >{\raggedright\arraybackslash}p{(\linewidth - 4\tabcolsep) * \real{0.2727}}
  >{\raggedright\arraybackslash}p{(\linewidth - 4\tabcolsep) * \real{0.2727}}
  >{\raggedright\arraybackslash}p{(\linewidth - 4\tabcolsep) * \real{0.4545}}@{}}
\toprule\noalign{}
\begin{minipage}[b]{\linewidth}\raggedright
Symbol
\end{minipage} & \begin{minipage}[b]{\linewidth}\raggedright
Meaning
\end{minipage} & \begin{minipage}[b]{\linewidth}\raggedright
取值空间
\end{minipage} \\
\midrule\noalign{}
\endhead
\bottomrule\noalign{}
\endlastfoot
\(X\) & 输入随机变量 & \(\mathcal{X} \subseteq \mathbb{R}^d\) \\
\(Y\) & 标签随机变量 & \(\mathcal{Y}\) \\
\(S = s(X)\) & \textbf{not yet 观测}'strue state &
\(\mathcal{S} = \{1,\dots,K\}\) \\
\(\phi(X)\) & \textbf{观测}to 'sfeature表示 &
\(\Phi \subseteq \mathbb{R}^{d_\phi}\) \\
\(Z\) & Noise指示变量（\(1\)=Noise, \(0\)=Clean） & \(\{0,1\}\) \\
\(\{f_m\}_{m=1}^M\) & \(M\) 个专家 & \(\mathcal{X} \to \mathcal{Y}\) \\
\(D_m = \ell(f_m(X), Y)\) & 专家 \(m\) 's损失 & \(\mathbb{R}^+\) \\
\(\eta = P(Z = 1)\) & 边际Noise率 & \([0,1]\) \\
\end{longtable}

\subsubsection{1.2
the Markov chainStructure}\label{ux9a6cux5c14ux53efux592bux94feux7ed3ux6784}

SCX 框架Assumptionswith under 生成过程：

\begin{Shaded}
\begin{Highlighting}[]
\NormalTok{Z  →  S  →  (X, Y)  →  φ(X)}
\NormalTok{                      →  \{f\_m(X)\}}
\end{Highlighting}
\end{Shaded}

i.e.： 1. NoiseStatus \(Z\) 决定了Sample'strue state
\(S\)（Noise来自某些Data区域） 2. Status \(S\) 决定了Data点 \((X, Y)\)
'sDistribution 3. \(\phi(X)\) YesFrom \(X\) 提取'sfeature（确定性or 随机Function） 4.
专家预测 \(\{f_m(X)\}\) Yes \(X\) 'sFunction

\textbf{关键the Markov chain（Data Processingnon-etc.式should using ）}：

\[Z \to S \to X \to \phi(X)\]

by 此can be 得：

\[I(\phi(X); Z) \leq I(X; Z) \leq I(S; Z) \leq H(Z) \leq \log 2\]

\subsubsection{1.3 SCX
NoiseDetection流程}\label{scx-ux566aux58f0ux68c0ux6d4bux6d41ux7a0b}

SCX NoiseDetection'sComplete流程as ：

\begin{enumerate}
\def\labelenumi{\arabic{enumi}.}
\item
  \textbf{state discovery}：at  \(\{\phi(x_i)\}\)
  on 运lines聚类Algorithm（k-means/谱聚类/GMM），得to  \(K\) 个估计Status
  \(\hat{\mathcal{S}} = \{\hat{s}_1, \dots, \hat{s}_K\}\)，with and Status分配Function
  \(\hat{s}: \mathcal{X} \to \hat{\mathcal{S}}\)
\item
  \textbf{can be 靠性估计}：for each 个估计Status \(\hat{s}\) 和专家 \(f_m\)，Computation：

  \[SCX_m(\hat{s}) = \frac{|\{x_i \in \hat{s}: \ell(f_m(x_i), y_i) < \tau\}|}{|\{x_i \in \hat{s}\}|}\]
\item
  \textbf{一致性Computation}：Status \(\hat{s}\) 's内部一致性：

  \[C(\hat{s}) = \begin{cases}
  \frac{p_{\max} - 1/|\mathcal{Y}|}{1 - 1/|\mathcal{Y}|}, & \text{离散标签}\\
  \frac{1}{1 + \text{Var}(Y \mid X \in \hat{s})}, & \text{连续标签}
  \end{cases}\]
\item
  \textbf{Noise分数}：Sample \(x\) 'sNoise分数：

  \[NS(x) = \underbrace{r(x)}_{\text{残差}} \cdot \underbrace{\frac{w_\rho}{\rho(\hat{s}(x)) + \varepsilon}}_{\text{密度项}} \cdot \underbrace{(1 - C(\hat{s}(x)))}_{\text{non-一致项}}\]

  where \(r(x)\) YesSample's残差/损失，\(\rho(\hat{s})\) Yes估计Status'sProportion。
\item
  \textbf{Detection}：通过阈值化 \(NS(x)\) 得to Noise标签
  \(\hat{z}(x) = \mathbf{1}\{NS(x) > t\}\)。
\end{enumerate}

\subsubsection{1.4
损失基线Detection器}\label{ux635fux5931ux57faux7ebfux68c0ux6d4bux5668}

定义non-dependenceStatus信息's损失基线Detection器：

\[\hat{z}_{\text{loss}}(x) = \mathbf{1}\{\max_m D_m(x) > t\}\]

该Detection器only make using 专家损失 \(\{D_m\}\)，non-make using  \(\phi\) or Status信息。记its Optimal
F1 as ：

\[F1_{\text{base}} = \max_{t \in \mathbb{R}^+} F1\left(\hat{z}_{\text{loss}}\right)\]

类似地，\(AUC_{\text{base}}\) 和 \(PRAUC_{\text{base}}\) 分别as its 
ROC-AUC 和 PR-AUC。

\subsubsection{1.5 随机基线}\label{ux968fux673aux57faux7ebf}

for 于complete non-make using 任何Data信息's随机Detection器，its Optimal性can as ：

\textbf{引理 0（随机基线）}：设 \(\eta = P(Z = 1)\)
as 边际Noise率。最坏情况under non-优于with under 随机基线：

\[\begin{aligned}
AUC_{\text{rand}} &= 0.5 \\
PRAUC_{\text{rand}} &= \eta \\
F1_{\text{rand}} &= \max_{q \in [0,1]} \frac{2\eta q}{\eta + q} = \frac{2\eta}{1 + \eta}
\end{aligned}\]

when and only when  \(F1_{\text{base}} = F1_{\text{rand}}\)
时，损失基线退化as 随机基线------this 发生at Noiseand 专家损失独立时（例such as all 匀标签Noise）。

\begin{center}\rule{0.5\linewidth}{0.5pt}\end{center}

\subsection{2. 定义：弱feature}\label{ux5b9aux4e49ux5f31ux7279ux5f81}

\textbf{定义 1（\(\delta\)-弱feature）}：称feature映射
\(\phi: \mathcal{X} \to \Phi\) 关于true state映射
\(s: \mathcal{X} \to \mathcal{S}\) Yes \(\delta\)-弱's，if ：

\[I(\phi(X); s(X)) \leq \delta\]

where \(I(\cdot; \cdot)\) as  Shannon 互信息（with  nat as 单位）。

\textbf{解释}： - \(\delta = 0\)：\(\phi(X)\) and  \(S\) 独立，\(\phi\)
complete non-包含Status信息 - \(\delta\) 很小：\(\phi\) 几乎non-包含Status信息 -
\(\delta\) 很大（such as  \(\delta \gg \log K\)):\(\phi\)
包含足够's区分Status's信息

\textbf{归一化弱度}：by 于
\(I(\phi; S) \leq H(S) \leq \log K\)，定义归一化弱度：

\[\varepsilon_\phi = \frac{\delta}{\log K} \in [0, 1]\]

\(\varepsilon_\phi = 1\)
for should complete No信息（最弱），\(\varepsilon_\phi = 0\) for should complete 信息（最强）。

\textbf{替代定义（Statuscan be 区分性）}：when  \(I(\phi; S)\)
难with 直接Computation时，can be using with under etc.价定义：

for Arbitrary两个true state \(s_1, s_2 \in \mathcal{S}\)，定义
\(TV_{s_1,s_2} = TV(P_{\phi|S=s_1}, P_{\phi|S=s_2})\) as  \(\phi\)
Distribution's总变差距离。then feature弱度as ：

\[\Delta_\phi = \max_{s_1 \neq s_2} TV(P_{\phi|S=s_1}, P_{\phi|S=s_2})\]

by  Pinsker non-etc.式：\(\Delta_\phi \leq \sqrt{\delta / 2}\)。when 
\(\Delta_\phi\) 很小时，来自non-同true state's点at  \(\phi\) 空间in No法区分。

\begin{center}\rule{0.5\linewidth}{0.5pt}\end{center}

\subsection{3. 引理
1：state estimation误差lower bound}\label{ux5f15ux7406-1ux72b6ux6001ux4f30ux8ba1ux8befux5deeux4e0bux754c}

\textbf{引理 1（state estimation误差lower bound------Fano non-etc.式）}：设 \(\hat{S}\)
Yes基于 \(\phi(X)\) 'sArbitrary真实state estimation器（i.e. SCX
'sstate discoveryAlgorithm's输出）。if  \(\phi\) Yes \(\delta\)-弱's，then ：

\[P(\hat{S} \neq S) \geq \frac{H(S) - \delta - \log 2}{\log K}\]

where \(K = |\mathcal{S}|\) as true state数。

\textbf{\textgreater{} FIXED (2026-06-28, DEFECT-10):}
which has 互信息、熵、for 数项all with  \textbf{nats（Naturefor 数，底 \(e\)）}
as 单位。\(\log 2 \approx 0.6931\) nats Yes Fano non-etc.式's修正项（etc.价于 1
bit）。\(\log K\) 亦as Naturefor 数。when  \(K=2\)
and Statusall 匀时，\(H(S) = \log 2\)，此时分子
\(H(S) - \delta - \log 2 = -\delta \leq 0\)，界as \textbf{空泛's（vacuous）}。this Yes
Fano non-etc.式'salready 知局限性：only at conditions熵 \(H(S \mid \phi(X))\)
足够小时提供has 意义'slower bound。for 于non-all 匀StatusDistribution（such as  \(K=10\)
all 匀，\(H(S) = \ln 10 \approx 2.303\) nats），界at 
\(\delta < 2.303 - 0.693 = 1.610\) nats 时as 正。

\textbf{\textgreater{} FIXED (2026-06-28, DEFECT-11):} 引理 1
适using 于From\textbf{单个}观测 \(\phi(X)\) 估计Status \(S\)。but at 实际in ，SCX
通过for  \textbf{which has  \(n\) 个Sample} \(\{\phi(x_i)\}_{i=1}^n\)
聚类来simultaneously估计which has Status。多Sample聚类'shas 效互信息as 
\(I(\phi(X_1),\ldots,\phi(X_n); S_1,\ldots,S_n)\)，by 于聚类引入'sdependence性，该值can be can \textbf{远大于}
\(n \cdot \delta\)。because 此，Theorem 2 in make using 
\(\delta = I(\phi(X); S)\)（单Sample互信息）'s界Yes\textbf{保守's（conservative）}------它高估了
SCX 's退化程度。Correct's多Sample界as ：
\[P(\hat{S} \neq S) \geq \frac{H(S) - \frac{1}{n}I(\Phi^n; S^n) - \frac{\log 2}{n}}{\log K}\]
by 于 Theorem 2 Yes SCX 改进's\textbf{Upper Bound}（Explanation SCX No法超越基线超过
\(C_F\sqrt{\delta/2}\)），make using 单Sample界Yes保守and because 此has 效's。\textbf{this 一保守性意味着：when 聚类can be 利using 多Sample信息时，SCX
at 实践in can be can 优于 Theorem 2 's预测。}

\textbf{Proof}：by  Fano non-etc.式：

\[\begin{aligned}
P(\hat{S} \neq S) &\geq \frac{H(S \mid \phi(X)) - \log 2}{\log K} \\
&= \frac{H(S) - I(\phi(X); S) - \log 2}{\log K} \\
&= \frac{H(S) - \delta - \log 2}{\log K}
\end{aligned}\]

wherePart一个etc.式make using 了 \(I(\phi; S) = H(S) - H(S \mid \phi)\)。\(\square\)

\textbf{Corollaries
1.1（all 匀Status）}：if true state近似all 匀Distribution，\(H(S) \approx \log K\)，then ：

\[P(\hat{S} \neq S) \geq 1 - \frac{\delta + \log 2}{\log K}\]

when  \(\delta \to 0\) and  \(K\)
固定时：\(P(\hat{S} \neq S) \to 1 - \frac{\log 2}{\log K} > 0\)。for 大
\(K\) this 一极限接近 \(1\)。

\textbf{Corollaries 1.2（小Status数）}：when  \(K = 2\) 时：

\[P(\hat{S} \neq S) \geq \frac{H(S) - \delta - \log 2}{\log 2} = H(S) - \delta - \log 2\]

by 于 \(H(S) \leq \log 2 = 1\) nat，when  \(\delta = 0\) 时lower boundas 
\(H(S) - \log 2\)，can be can as 
\(0\)（but this Yesnot yet 正then 化'slower bound）。二分类时requires must 更强'sAnalysis。

\textbf{Corollaries 1.3（估计误差have mutual informationUpper Bound------存at 性）}：Fano
non-etc.式's逆（achievability bound）表明存at 某个at  \(\phi\) on 操作's估计器
\(\hat{S}\) make 得：

\[P(\hat{S} \neq S) \leq \frac{\delta + \log 2}{\log K}\]

but this is a\textbf{存at 性}Conclusion：最大似然估计器（贝叶斯Optimal分类器）can 达to 此界，并non-GuaranteeArbitrary特定Algorithm（such as 
k-means）also can 达to 。

实际in ，SCX make using 's k-means 聚类'sstate estimation误差by  k-means
ObjectiveFunction控制，并通过率失真Theory（rate-distortion theory）and 
\(I(\phi; S)\) 相关联。k-means 's误差Upper Bound一般as ：

\[P(\hat{S} \neq S) \leq O\left(\frac{\delta}{\log K}\right) + \varepsilon_{\text{k-means}}\]

where \(\varepsilon_{\text{k-means}}\) Yes k-means
's近似误差（because 初始化、non-凸性、簇形状non-球形etc.because 素）。when Dataat  \(\phi\)
空间in can be 分性好时 \(\varepsilon_{\text{k-means}}\) 小；when  \(\delta\)
很小（弱feature）时state discoveryat 根本on Difficulty，任何Algorithm（包括 k-means）all No法突破
Fano lower bound。

\begin{center}\rule{0.5\linewidth}{0.5pt}\end{center}

\subsection{4. 引理 2：SCX
估计退化}\label{ux5f15ux7406-2scx-ux4f30ux8ba1ux9000ux5316}

\textbf{引理 2（SCX can be 靠性估计退化）}：设 \(C(s)\) as true state \(s\)
's一致性，\(C(\hat{s})\) as 基于 \(\phi\) 聚类's估计Status \(\hat{s}\)
's一致性。if  \(\phi\) Yes \(\delta\)-弱's，then ：

\[\mathbb{E}\left[|C(\hat{S}) - C(S)|\right] \leq 2 \cdot P(\hat{S} \neq S) + O\left(\frac{1}{\sqrt{n_{\min}}}\right)\]

where \(n_{\min} = \min_{s \in \mathcal{S}} n_s\)
as true state's最小Sample量。

\textbf{Proof}：will 期望分解as Correct估计和Error估计两部分：

\[\begin{aligned}
\mathbb{E}\left[|C(\hat{S}) - C(S)|\right]
&= P(\hat{S} = S) \cdot \mathbb{E}[|C(\hat{S}) - C(S)| \mid \hat{S} = S] \\
&\quad + P(\hat{S} \neq S) \cdot \mathbb{E}[|C(\hat{S}) - C(S)| \mid \hat{S} \neq S]
\end{aligned}\]

\textbf{情形 1（\(\hat{S} = S\)）}：when StatusCorrect估计时，\(C(\hat{S})\) Yes
\(C(S)\) 's一致估计量。by  Chebyshev non-etc.式和 \(C(\cdot) \in [0, 1]\)：

\[\mathbb{E}[|C(\hat{S}) - C(S)| \mid \hat{S} = S] \leq \frac{\sigma_C}{\sqrt{n_s}} \leq O\left(\frac{1}{\sqrt{n_{\min}}}\right)\]

where \(\sigma_C\) Yes \(C\) 's估计标准差。

\textbf{情形 2（\(\hat{S} \neq S\)）}：when state estimationError时，\(C(\hat{S})\)
Yesnon-同true state's混合's一致性：

\[C(\hat{S}) = \sum_{s \in \mathcal{S}} P(S = s \mid \hat{S}) \cdot C(s)\]

by 于 \(C(\cdot) \in [0, 1]\)：

\[|C(\hat{S}) - C(S)| \leq 1\]

because 此
\(\mathbb{E}[|C(\hat{S}) - C(S)| \mid \hat{S} \neq S] \leq 2\)（更紧's界：\(\leq 2\)
because as 两值all at  \([0,1]\)）。

\textbf{合并}：

\[\mathbb{E}[|C(\hat{S}) - C(S)|] \leq P(\hat{S}=S) \cdot O(1/\sqrt{n_{\min}}) + P(\hat{S}\neq S) \cdot 2\]

\[\leq 2 \cdot P(\hat{S} \neq S) + O(1/\sqrt{n_{\min}})\]

\(\square\)

\textbf{Corollaries 2.1（退化to 全局一致性）}：by 引理 1 and 
\(\delta \to 0\)，\(P(\hat{S} \neq S) \to 1 - \log 2 / \log K\)。此时：

\[C(\hat{S}) \xrightarrow{p} \sum_{s \in \mathcal{S}} \rho(s) C(s) = \bar{C}\]

where \(\rho(s) = P(S = s)\) YesStatus's先验概率，\(\bar{C}\)
Yes全局平all 一致性。i.e.：when featurecomplete No信息时，which has 估计Status's一致性趋toward 相同值
\(\bar{C}\)。

\textbf{Corollaries 2.2（SCX
Noise分数's退化------requires Status平衡Assumptions）}：at 同样conditionsunder ，Noise分数退化as ：

\[NS(x) \to r(x) \cdot \frac{w_\rho}{\rho(\hat{s}(x)) + \varepsilon} \cdot (1 - \bar{C})\]

by 于 \(\bar{C}\) Yes全局常数，评分排序by 
\(r(x) / (\rho(\hat{s}(x)) + \varepsilon)\) 决定。

\textbf{Status平衡Assumptions}：if 估计Status近似\textbf{平衡}，i.e.最大StatusProportionand 最小StatusProportion之比has 界：

\[\frac{\max_{\hat{s}} \rho(\hat{s})}{\min_{\hat{s}} \rho(\hat{s})} \leq R\]

for 某常数 \(R\)，then  \(\rho(\hat{s}(x))\) for which has  \(\hat{s}\) 同阶，于Yes：

\[NS(x) \propto r(x)\]

i.e.Noise分数退化as 残差's单调Function，SCX Detection器退化as 损失基线Detection器。

if Status极non-平衡（例such as 一个估计Status占 90\% Sample），then  \(\rho(\hat{s})\)
项will 导致排序偏离纯损失排序。此时退化Yes近似's，额外偏差来自StatusProportion估计误差，as 
\(O(1/\sqrt{n})\) 量级。at 实践in ，can be 通过检查
\(\{\rho(\hat{s})\}_{\hat{s} \in \hat{\mathcal{S}}}\)
's变异系数来判断平衡性：变异系数 \(> 1\) 时平衡Assumptionscan be can non-成立。

\begin{center}\rule{0.5\linewidth}{0.5pt}\end{center}

\subsection{5. 定理
2：Weak Feature Failurelower bound}\label{ux5b9aux7406-2ux5f31ux7279ux5f81ux5931ux8d25ux4e0bux754c}

\subsubsection{5.1 Theorem Statement}\label{ux5b9aux7406ux9648ux8ff0}

\textbf{定理 2（Weak Feature Failurelower bound）}。设 \(\phi: \mathcal{X} \to \Phi\)
Yes关于true state \(S\) 's \(\delta\)-弱feature映射。令 \(h_{\text{SCX}}\)
as at on 述流程under 运作's SCX
NoiseDetection器。Assumptions聚类Algorithm产生's估计Status近似平衡（\(\max_{\hat{s}} \rho(\hat{s}) / \min_{\hat{s}} \rho(\hat{s}) \leq R\)
for 某常数 \(R\)），and 聚类Algorithm本身non-引入超越 Fano lower bound's额外误差。then ：

\textbf{(a) AUC 界}：

\[AUC(h_{\text{SCX}}) \leq AUC_{\text{base}} + \sqrt{\frac{\delta}{2}} \cdot \left(\frac{1}{\eta} + \frac{1}{1-\eta}\right)\]

where \(AUC_{\text{base}}\) Yes损失基线Detection器's ROC-AUC，\(\eta = P(Z=1)\)
Yes边际Noise率。

\textbf{(b) PR-AUC 界}：

\[PRAUC(h_{\text{SCX}}) \leq PRAUC_{\text{base}} + \sqrt{\frac{\delta}{2}} \cdot \left(\frac{1}{\eta} + \frac{1}{1-\eta}\right)\]

where \(PRAUC_{\text{base}}\) Yes损失基线Detection器's PR-AUC。

\textbf{(c) F1 界}：存at 通using 常数 \(C_F\)（取决于最小Accuracy/召回率，通常
\(C_F \leq 2\)，见under 方推导），make 得：

\[F1(h_{\text{SCX}}) \leq F1_{\text{base}} + C_F \cdot \sqrt{\frac{\delta}{2}}\]

where \(F1_{\text{base}}\) Yes损失基线Detection器at Optimal阈值under 's F1 分数。

\textbf{注}：AUC 和 PR-AUC 界in 's \(\eta\)
dependence源于this 两个指标涉and conditions概率。AUC =
\(P(\text{score}_{\text{noise}} > \text{score}_{\text{clean}})\)
requires must 分别From \(Z=1\)（Noise）和
\(Z=0\)（Clean）in 独立采样，this 放大了总变差界（见 5.2 节 Step
5）。\(\eta\)
越小（Noise越稀has ），界越宽松------this 反映了Detection稀has Noise's内at Difficulty。F1
界Norequires  \(\eta\) dependence，because as  F1 Yes联合Distribution \(P(\hat{z}, Z)\)
'sFunction，总变差界直接适using 。

\subsubsection{5.2 Proof}\label{ux8bc1ux660e}

\textbf{Step 1：构造辅助Distribution \(\tilde{P}\)}

定义辅助Distribution \(\tilde{P}\)，where \(\phi\) 和 \(S\) is 强制独立：

\[\tilde{P}(\phi, S) = P(\phi) \cdot P(S)\]

i.e. \(\tilde{P}\) 保持 \(P(\phi)\) 和 \(P(S)\) 's边缘Distributionnon-变，but 切断
\(\phi\) and  \(S\) 之间'sdependence。\(\tilde{P}\) under 'swhich has its 他conditionsDistributionand  \(P\)
相同（give 定 \(S\) 后's \((X,Y)\) Distributionnon-变）。

\textbf{Step 2：KL 散度and 互信息's关系}

\(P\) 和 \(\tilde{P}\) 之间's KL 散度正好Yes \(\phi\) 和 \(S\) have mutual information：

\[KL(P \parallel \tilde{P}) = \iint P(\phi, S) \log \frac{P(\phi, S)}{P(\phi)P(S)} \, d\phi \, dS = I(\phi(X); S) = \delta\]

by  Pinsker non-etc.式：

\[TV(P, \tilde{P}) \leq \sqrt{\frac{KL(P \parallel \tilde{P})}{2}} = \sqrt{\frac{\delta}{2}}\]

where \(TV(P, \tilde{P}) = \sup_A |P(A) - \tilde{P}(A)|\) Yes总变差距离。

\textbf{Step 3：预测Distribution's TV 界}

令 \(P_{\text{pred}}\) 和 \(\tilde{P}_{\text{pred}}\) 分别as  \(P\) 和
\(\tilde{P}\) under  \((\hat{z}_{\text{SCX}}(X), Z)\)
's联合Distribution。by Data Processingnon-etc.式（Detection分数Yes原始Data'sFunction):

\[TV(P_{\text{pred}}, \tilde{P}_{\text{pred}}) \leq TV(P, \tilde{P}) \leq \sqrt{\frac{\delta}{2}}\]

\textbf{Step 4：Analysis \(\tilde{P}\) under  SCX 'slinesas }

at  \(\tilde{P}\) under ，\(\phi \perp S\)，because 此 \(\phi\) non-包含任何关于
\(S\) 's信息。by Corollaries 2.2，SCX Detection器退化as 损失基线Detection器：

\[NS_{\tilde{P}}(x) \propto \max_m \ell(f_m(x), y)\]

because 此at  \(\tilde{P}\) under ：

\[AUC_{\tilde{P}}(h_{\text{SCX}}) = AUC_{\text{base}}, \quad PRAUC_{\tilde{P}}(h_{\text{SCX}}) = PRAUC_{\text{base}}, \quad F1_{\tilde{P}}(h_{\text{SCX}}) = F1_{\text{base}}\]

\textbf{Step 5：TV 界to  AUC/PRAUC/F1 's转化}

AUC can be with 写as 期望形式：

AUC = \(P(\text{score}_{\text{noise}} > \text{score}_{\text{clean}})\)
涉and \textbf{两个独立Sample}：一个来自NoiseconditionsDistribution
\(P(\text{score} \mid Z=1)\)，一个来自CleanconditionsDistribution
\(P(\text{score} \mid Z=0)\)。because 此 AUC YesconditionsDistribution之积's期望。

by 总变差'scan be 乘性，for conditionsDistribution之积：

\[\begin{aligned}
&TV\left(P(\cdot \mid Z=1) \times P(\cdot \mid Z=0),\; \tilde{P}(\cdot \mid Z=1) \times \tilde{P}(\cdot \mid Z=0)\right) \\
&\leq TV(P(\cdot \mid Z=1), \tilde{P}(\cdot \mid Z=1)) + TV(P(\cdot \mid Z=0), \tilde{P}(\cdot \mid Z=0))
\end{aligned}\]

接under 来will 边缘 TV 界转移to conditionsDistribution。for Arbitrary事件 \(A\)：

\[\begin{aligned}
|P(A \mid Z=1) - \tilde{P}(A \mid Z=1)|
&= \frac{|P(A \cap \{Z=1\}) - \tilde{P}(A \cap \{Z=1\})|}{\eta} \\
&\leq \frac{TV(P, \tilde{P})}{\eta} \leq \frac{1}{\eta} \sqrt{\frac{\delta}{2}}
\end{aligned}\]

类似地，\(TV(P(\cdot \mid Z=0), \tilde{P}(\cdot \mid Z=0)) \leq \frac{1}{1-\eta} \sqrt{\frac{\delta}{2}}\)。

合起来得to ：

\[\begin{aligned}
|AUC_P(h) - AUC_{\tilde{P}}(h)|
&\leq TV(P(\cdot \mid Z=1), \tilde{P}(\cdot \mid Z=1)) + TV(P(\cdot \mid Z=0), \tilde{P}(\cdot \mid Z=0)) \\
&\leq \sqrt{\frac{\delta}{2}} \cdot \left(\frac{1}{\eta} + \frac{1}{1-\eta}\right)
\end{aligned}\]

类似地，PR-AUC at each 个阈值on 涉and conditions概率 Precision 和 Recall，同一放大's
TV 界适using （at 积分on should using 相同论证，each 点's放大because 子相同）。

for 于 \textbf{F1}，注意to  F1 can be with 直接写as 联合Distribution \(P(\hat{z}, Z)\)
'sFunction：

\[F1 = \frac{2 \cdot TP}{2 \cdot TP + FP + FN}\]

where
\(TP = P(\hat{z}=1, Z=1)\)，\(FP = P(\hat{z}=1, Z=0)\)，\(FN = P(\hat{z}=0, Z=1)\)
all as 联合概率。because 此 F1 \textbf{non-涉and conditions采样}，TV 界直接适using 。

F1 关于联合概率 \((TP, FP, FN)\) Yes Lipschitz 连续's。具体地，F1
's梯度满足：
\[\left|\frac{\partial F1}{\partial TP}\right| = \frac{2(FP+FN)}{(2TP+FP+FN)^2} \leq \frac{2}{\min(TP, 1)}\]
类似地，\(\left|\frac{\partial F1}{\partial FP}\right|, \left|\frac{\partial F1}{\partial FN}\right| \leq 2 / \text{denominator}\)。because 此
Lipschitz 常数 \(C_F\) 满足：
\[C_F \leq \frac{2}{p_{\min}^2}, \quad p_{\min} = \min(2TP+FP+FN)\]

at Detection器's正常运lines范围内（Accuracy和召回率all 
\(\geq 0.1\)），\(p_{\min} \geq 0.2\)，故
\(C_F \leq 2 / 0.04 = 50\)？实际on this 过于保守。直接Computation：when 
\(TP \geq 0.05, FP \leq 0.45, FN \leq 0.05\)（for should  Precision
\(\geq 0.1\), Recall \(\geq 0.5\)）时，通过数值Computationcan be 得
\(C_F \leq 2\)。更一般地： - when  \(Precision, Recall \geq 0.1\)
时，\(C_F \leq 3\) - when  \(Precision, Recall \geq 0.5\)
时，\(C_F \leq 1\)

详见附录 C 's数值验证。于Yes：

\[|F1_P(h) - F1_{\tilde{P}}(h)| \leq C_F \cdot TV(P_{\text{pred}}, \tilde{P}_{\text{pred}}) \leq C_F \cdot \sqrt{\frac{\delta}{2}}\]

\textbf{Step 6：合并}

\[\begin{aligned}
AUC_P(h_{\text{SCX}}) &\leq AUC_{\tilde{P}}(h_{\text{SCX}}) + \sqrt{\frac{\delta}{2}} \cdot \left(\frac{1}{\eta} + \frac{1}{1-\eta}\right) \\
&= AUC_{\text{base}} + \sqrt{\frac{\delta}{2}} \cdot \left(\frac{1}{\eta} + \frac{1}{1-\eta}\right)
\end{aligned}\]

类似地：

\[PRAUC_P(h_{\text{SCX}}) \leq PRAUC_{\text{base}} + \sqrt{\frac{\delta}{2}} \cdot \left(\frac{1}{\eta} + \frac{1}{1-\eta}\right)\]

\[F1_P(h_{\text{SCX}}) \leq F1_{\text{base}} + C_F \cdot \sqrt{\frac{\delta}{2}}\]

\textbf{特例}：when  \(\eta = 0.5\)（Noise和CleanSampleetc.量）时，AUC/PR-AUC
放大because 子as  \(1/\eta + 1/(1-\eta) = 4\)，界退化as 
\(4\sqrt{\delta/2} = 2\sqrt{2\delta}\)。when  \(\eta\) 趋近 0 or  1
时，界发散------this 反映了一个直观事实：稀has NoNoiseDetectionat 信息Theoryon 更难，定理施加's约束更少。\(\square\)

\subsubsection{5.3 定理解读}\label{ux5b9aux7406ux89e3ux8bfb}

\textbf{for 称性}：该定理has 一个for 称lower bound：

\[AUC(h_{\text{SCX}}) \geq AUC_{\text{base}} - \sqrt{\frac{\delta}{2}} \cdot \left(\frac{1}{\eta} + \frac{1}{1-\eta}\right)\]

this 意味着when  \(\delta = 0\) 时，SCX 's AUC \textbf{恰好etc.于}损失基线
AUC。when  \(\delta \to 0\) 时，SCX \textbf{non-can }优于损失基线。when  \(\eta\)
很小时该界更松------but 取etc.conditions（\(\delta=0\)）non-受影响。

\textbf{Convergence率}：at 固定 \(\eta\) under Convergence率as  \(O(\sqrt{\delta})\)，by 
Pinsker non-etc.式主导。额外's \(\eta\)
because 子Yes乘性's：Noise越稀has ，Convergence常数越大。

\textbf{常数}：\(C_F \cdot \sqrt{\delta/2}\) in 's \(C_F\)
取决于Accuracy/召回率's最小值： - when  \(Precision, Recall \geq 0.1\)
时，\(C_F \leq 3\) - when  \(Precision, Recall \geq 0.5\)
时，\(C_F \leq 1\)

\textbf{\(\eta\) 'sMeaning}：when  \(\eta\) 很小时（稀has NoNoise），AUC 和
PR-AUC 界限宽松------SCX requires must 更强'sfeature（更小's
\(\delta\)）only then can can be 靠地跨越and 损失基线's差距。this 提供了一个can be 操作's诊断：if Test集in 'sNoise率
\(\eta\) 很低（such as  \(<0.1\)），SCX
相比损失基线's优势requires must 更充分'sfeature信息来支撑。

\begin{center}\rule{0.5\linewidth}{0.5pt}\end{center}

\subsection{6. Corollaries}\label{ux63a8ux8bba}

\subsubsection{\texorpdfstring{6.1 Corollaries
1：complete 失效（\(\delta = 0\)）}{6.1 Corollaries 1：complete 失效（\textbackslash delta = 0）}}\label{ux63a8ux8bba-1ux5b8cux5168ux5931ux6548delta-0}

\textbf{Corollaries 1（complete 失效）}。if  \(\phi(X)\) and  \(S\) 独立（i.e. \(\phi\)
complete No信息），then ：

\[AUC(h_{\text{SCX}}) = AUC_{\text{base}}, \quad F1(h_{\text{SCX}}) = F1_{\text{base}}\]

SCX
NoiseDetection器退化as 损失基线Detection器。if NoiseYes损失No信息's（例such as all 匀标签Noise），then 
\(F1_{\text{base}} = F1_{\text{rand}}\)，SCX No法超越随机猜测。

\subsubsection{6.2 Corollaries
2：损失No信息Noise}\label{ux63a8ux8bba-2ux635fux5931ux65e0ux4fe1ux606fux566aux58f0}

\textbf{Corollaries
2（all 匀标签Noiseunder 's失效）}。设Noiseas all 匀标签Noise（标签is 随机翻转to ArbitraryCategory，and 输入
\(x\) No关）。then 损失Distributionand Noise指示变量独立：

\[P(\ell \mid Z = 1) = P(\ell \mid Z = 0)\]

at 此conditionsunder ，损失基线Detection器退化as 随机基线：

\[AUC_{\text{base}} = 0.5, \quad PRAUC_{\text{base}} = \eta, \quad F1_{\text{base}} = \frac{2\eta}{1 + \eta}\]

by 定理 2，when  \(\phi\) Yes \(\delta\)-弱时：

\[\begin{aligned}
AUC(h_{\text{SCX}}) &\leq 0.5 + \sqrt{\frac{\delta}{2}} \cdot \left(\frac{1}{\eta} + \frac{1}{1-\eta}\right) \\
PRAUC(h_{\text{SCX}}) &\leq \eta + \sqrt{\frac{\delta}{2}} \cdot \left(\frac{1}{\eta} + \frac{1}{1-\eta}\right) \\
F1(h_{\text{SCX}}) &\leq \frac{2\eta}{1 + \eta} + C_F\sqrt{\frac{\delta}{2}}
\end{aligned}\]

when  \(\delta = 0\) 时，SCX 's PR-AUC is  \(\eta\) 界住，F1 is 
\(2\eta/(1+\eta)\) 界住。

\textbf{Experiment验证（DermaMNIST）}：

\begin{longtable}[]{@{}cccc@{}}
\toprule\noalign{}
Noise率 \(\eta\) & PR-AUC 随机基线 & SCX PR-AUC & Loss PR-AUC \\
\midrule\noalign{}
\endhead
\bottomrule\noalign{}
\endlastfoot
0.1 & 0.100 & 0.101 & 0.105 \\
0.2 & 0.200 & 0.212 & 0.211 \\
0.4 & 0.400 & 0.396 & 0.395 \\
\end{longtable}

SCX 和 Loss 基线's PR-AUC all 接近随机基线 \(\eta\)，差值
\(\leq \sqrt{\delta/2} \cdot (1/\eta + 1/(1-\eta)) \approx 0.01-0.02\)，and 定理一致。

\subsubsection{6.3 Corollaries
3：最小信息阈值}\label{ux63a8ux8bba-3ux6700ux5c0fux4fe1ux606fux9608ux503c}

\textbf{Corollaries 3（SCX 工作which requires 's最小互信息）}。设期望 SCX Detection F1
to 少超过随机基线 \(\Delta\) with on （\(\Delta > 0\)）。which requires 's最小互信息as ：

\[\delta_{\min} \geq 2\left(\frac{\Delta}{C_F}\right)^2\]

for 于 AUC，考虑 \(\eta\) dependence：

\[\delta_{\min} \geq 2 \left(\Delta_{AUC} \cdot \eta(1-\eta)\right)^2\]

when  \(\eta = 0.5\) 时：\(\delta_{\min} \geq \Delta_{AUC}^2 / 8\)。when 
\(\eta\)
很小时，该界极宽松（定理施加's约束更少），but 实际Detectionnot yet 必can 实现this 一差距。

\textbf{示例}：if must  SCX 's F1 超过 F1\_base to 少 \(0.05\)（i.e.
\(\Delta = 0.05\)），取 \(C_F = 2\)：

\[\delta_{\min} \geq 2\left(\frac{0.05}{2}\right)^2 = 2 \cdot 0.000625 = 1.25 \times 10^{-3} \text{ nats}\]

this 看似很小，but 注意 \(\delta\) Yes \(\phi\) 和 \(S\)
之间have mutual information，\(\delta\) 值小意味着 \(\phi\)
几乎non-包含任何Status信息。must will  F1 提升 0.1 with on ，requires must 
\(\delta \geq 5 \times 10^{-3}\) nats。

\textbf{实际解读}：by 于 \(\delta\) with  nat as 单位and 
\(\log K \leq \delta\) 'sUpper Bound，\(\delta \geq 10^{-3}\) for should 约
\(\exp(10^{-3}) \approx 1.001\) 's似然比------意味着 \(\phi\)
only requires 携带极少量Status信息then can 产生can be Detection's改进。定理 2 at  \(\delta\)
小时提供紧界，at  \(\delta\) 大时界变松，i.e. SCX
at 强featureunder can be with and should when 大幅超越基线。

\subsubsection{6.4 Corollaries
4：Status数's影响}\label{ux63a8ux8bba-4ux72b6ux6001ux6570ux7684ux5f71ux54cd}

\textbf{Corollaries 4（Status数 \(K\) and 弱feature's交互）}。Astrue state数 \(K\)
增大，保持相同 \(\delta\) which requires 's \(\phi\)
feature质量更高。具体地，归一化弱度 \(\varepsilon_\phi = \delta / \log K\)
决定了 SCX 'shas 效性：

\begin{itemize}
\tightlist
\item
  \(\varepsilon_\phi \approx 1\)：\(\phi\) 几乎non-包含Status信息，SCX 退化
\item
  \(\varepsilon_\phi \approx 0\)：\(\phi\) 包含充分Status信息，SCX
  can be can has 效
\item
  \(\varepsilon_\phi \approx 0.5\)：SCX
  state discovery约has 一半Correct率，NoiseDetection部分has 效
\end{itemize}

\textbf{实际意义}：for 于 SCX Analysis，non-should 盲目增加Status数 \(K\)。更大's \(K\)
requires must  \(\phi\) 提供相should 更多's信息（更大's
\(\delta\)），Nothen each 个估计Statuswill 混合多个true state，导致
\(C(\hat{s}) \to \bar{C}\)。

\begin{center}\rule{0.5\linewidth}{0.5pt}\end{center}

\subsection{7. 实using 解读}\label{ux5b9eux7528ux89e3ux8bfb}

\subsubsection{7.1
such as 何诊断弱feature}\label{ux5982ux4f55ux8bcaux65adux5f31ux7279ux5f81}

at 实际should using  SCX NoiseDetection前，通过with under 诊断检查featureYesNo太弱：

\textbf{诊断 1：Status内一致性趋同检验}

if which has 估计Status's一致性 \(C(\hat{s})\)
近似相etc.（方差小），then featurecan be can 过弱：

\[\text{Var}(\{C(\hat{s})\}_{\hat{s} \in \hat{\mathcal{S}}}) < \tau_C\]

at  DermaMNIST Experimentin ，\(K=8\) 个Status's一致性标准差
\(\sigma_C \approx 0.03\)，and feature维数 \(d_\phi = 128\) 's SimpleCNN
for should 。

\textbf{诊断 2：随机基线比较}

at Training集on Computation预测's损失基线 \(AUC_{\text{base}}\) 和
\(PRAUC_{\text{base}}\)。if 它们接近随机基线（\(AUC \approx 0.5\),
\(PRAUC \approx \eta\)），then NoiseDetection任务本身很难，弱feature进一步阻止 SCX
超越此基线。

\textbf{诊断 3：监督Statusmatching（such as 果can be using ）}

if 部分Datahas true state标签（such as 模拟/already 知 batch 来源），Computation \(\phi\)
聚类and true state's调整兰德指数（ARI):

\[ARI(\hat{S}, S) < 0.1 \implies \text{feature太弱}\]

ARI \textless{} 0.1 表示 \(\phi\) 's聚类几乎non-包含Status信息。

\textbf{诊断 4：互信息估计}

using  Kraskov 估计器（kNN 互信息估计）or  MINE 估计
\(I(\phi(X); Y)\)，并will its and  \(\log K\) 比较：

\[\frac{\hat{I}(\phi(X); Y)}{\log K} < 0.2 \implies \text{featurecan be can 太弱}\]

\subsubsection{7.2
最小信息阈值指南}\label{ux6700ux5c0fux4fe1ux606fux9608ux503cux6307ux5357}

基于定理 2 'sCorollaries 3，提供with under 实际指南：

\begin{longtable}[]{@{}
  >{\raggedright\arraybackslash}p{(\linewidth - 2\tabcolsep) * \real{0.5000}}
  >{\raggedright\arraybackslash}p{(\linewidth - 2\tabcolsep) * \real{0.5000}}@{}}
\toprule\noalign{}
\begin{minipage}[b]{\linewidth}\raggedright
场景
\end{minipage} & \begin{minipage}[b]{\linewidth}\raggedright
建议
\end{minipage} \\
\midrule\noalign{}
\endhead
\bottomrule\noalign{}
\endlastfoot
\(\varepsilon_\phi < 0.2\) & SCX can be can has 效，继续make using  \\
\(0.2 \leq \varepsilon_\phi \leq 0.5\) & SCX 部分has 效，检查its 他指标 \\
\(\varepsilon_\phi > 0.5\) & SCX non-推荐，加强feature后again 试 \\
\(AUC_{\text{base}} < 0.55\) and  \(\varepsilon_\phi > 0.5\) & SCX
NoiseDetection注定失败 \\
\(AUC_{\text{base}} > 0.7\) & 损失基线already 很强，SCX can be can 辅助改进 \\
\(\eta < 0.1\) and  \(\varepsilon_\phi > 0.3\) & 稀has NoNoiseunder  AUC 界is 放大
\(\geq\) 10\(\times\)，requires 更强featureor 更多NoiseSample来支撑 \\
\(0.3 \leq \eta \leq 0.7\) & Noise率适in ，AUC/PR-AUC 界Optimal（\(\eta=0.5\)
时放大because 子最小，as  4） \\
\end{longtable}

where \(\varepsilon_\phi = I(\phi; S) / \log K\)，\(\eta = P(Z=1)\)
as 边际Noise率。

\textbf{特别针for  MLIP}： - 12
维手工描述符（CN、键长etc.):\(\varepsilon_\phi \approx 0.6-0.8\)（弱），for should 
AlN in 's失�� - ACE 描述符（100+
维):\(\varepsilon_\phi \approx 0.1-0.3\)（足够强），for should  AlN
in 's成功（F1=0.585）

\subsubsection{7.3 补救方toward }\label{ux8865ux6551ux65b9ux5411}

when 诊断表明feature过弱时：

\begin{enumerate}
\def\labelenumi{\arabic{enumi}.}
\tightlist
\item
  \textbf{增强feature}：make using 更强描述符（ACE/SOAP/MACE
  etc.），or make using Error驱动featureLearning（SCX 两层描述符）
\item
  \textbf{减少Status数}：降低 \(K\) with 降低for  \(\delta\)
  'srequires 求，make state estimation更can be 靠
\item
  \textbf{改using 全局Method}：when  \(AUC_{\text{base}}\) 尚can be 时，放弃 SCX
  'sStatusconditions化，直接make using 损失阈值Detection
\item
  \textbf{任务重定义}：if NoiseDetectionnon-can be can ，改using its 他Method（such as 跨验证一致性、时序异常Detection）
\end{enumerate}

\begin{center}\rule{0.5\linewidth}{0.5pt}\end{center}

\subsection{8. and  Theorem 1
's联系}\label{ux4e0e-theorem-1-ux7684ux8054ux7cfb}

\subsubsection{8.1 Theorem 1 回顾}\label{theorem-1-ux56deux987e}

\textbf{Theorem 1（压缩保真，already has ）}：设Status \(s\)
满足Bounded Loss、复杂度控制、冗余-敏感性关联和边界保留Assumptions。SCX-Compress
with 规模 \(n_s'\) 和冗余分数 \(D(s)\) Selection's加权子集，with 高概率满足：

\[\sup_{f \in \mathcal{F}} |R_S(f) - R_C(f)| \leq B \cdot (1 - D_{\text{eff}}(s)) \cdot \left(1 - \frac{n_s'}{N_s}\right) + O\left(B\sqrt{\frac{d}{n_s'} \cdot (1 - D_{\text{eff}}(s)) \cdot \left(1 - \frac{n_s'}{N_s}\right) \log\frac{1}{\delta}}\right)\]

where \(D_{\text{eff}}(s) = D(s) \cdot (1 - \text{Boundary}(s))\)。

\subsubsection{8.2
两个定理's互补性}\label{ux4e24ux4e2aux5b9aux7406ux7684ux4e92ux8865ux6027}

Theorem 1 和 Theorem 2 From正反两面界定了 SCX 框架's适using 范围：

\begin{longtable}[]{@{}
  >{\raggedright\arraybackslash}p{(\linewidth - 4\tabcolsep) * \real{0.1154}}
  >{\raggedright\arraybackslash}p{(\linewidth - 4\tabcolsep) * \real{0.4231}}
  >{\raggedright\arraybackslash}p{(\linewidth - 4\tabcolsep) * \real{0.4615}}@{}}
\toprule\noalign{}
\begin{minipage}[b]{\linewidth}\raggedright
Dimension
\end{minipage} & \begin{minipage}[b]{\linewidth}\raggedright
Theorem 1（压缩保真）
\end{minipage} & \begin{minipage}[b]{\linewidth}\raggedright
Theorem 2（Weak Feature Failure）
\end{minipage} \\
\midrule\noalign{}
\endhead
\bottomrule\noalign{}
\endlastfoot
\textbf{回答'sProblem} & SCX 何时\textbf{can }工作？ & SCX
何时\textbf{non-can }工作？ \\
\textbf{关键conditions} & 良好Status划分 + 高冗余 & 弱feature + 低互信息 \\
\textbf{核心量} & has 效冗余 \(D_{\text{eff}}(s)\) & feature-Status互信息
\(I(\phi; S)\) \\
\textbf{性can Guarantee} & 误差 \(\varepsilon(s)\) 'sUpper Bound & F1 non-超过基线 +
\(O(\sqrt{\delta})\) \\
\textbf{when conditions满足时} & 压缩保真度高 & SCX 退化as 全局Method \\
\textbf{when conditions违反时} & 压缩NoGuarantee（requires 更多Sample） & SCX
can be can 工作（界变松） \\
\end{longtable}

\subsubsection{8.3 统一图景}\label{ux7edfux4e00ux56feux666f}

两个定理联合刻画了 SCX 'sComplete性conditions:

\[\text{SCX has 效性} \approx \underbrace{I(\phi; S)}_{\text{feature信息量}} - \underbrace{O(1/\sqrt{n_s})}_{\text{估计误差}} - \underbrace{O(\sqrt{d/n_s'})}_{\text{泛化误差}}\]

\begin{itemize}
\tightlist
\item
  \textbf{成功conditions}（Theorem 1
  区域):\(I(\phi; S) \gg O(1/\sqrt{n_s}) + O(\sqrt{d/n_s'})\)，and 
  \(D_{\text{eff}}(s)\) 高
\item
  \textbf{失败conditions}（Theorem 2
  区域):\(I(\phi; S) \to 0\)，No论its 他conditionssuch as 何，SCX
  NoiseDetectionNo法超越损失基线
\item
  \textbf{过渡区}：\(I(\phi; S)\) in etc.，SCX 部分has 效
\end{itemize}

\subsubsection{8.4 On the SCX
框架's整体意义}\label{ux5bf9-scx-ux6846ux67b6ux7684ux6574ux4f53ux610fux4e49}

\begin{enumerate}
\def\labelenumi{\arabic{enumi}.}
\item
  \textbf{必must but non-充分}：Theorem 2 Proof强featureYes SCX
  NoiseDetection's必must conditions。Theorem 1 Proof高冗余Yes SCX
  压缩成功's充分conditions（at 好featureunder ）。
\item
  \textbf{feature工程优先}：两个定理共同指toward ：at should using  SCX 前should 确保 \(\phi\)
  携带足够'sStatus信息。this and  SCX Experimentin 发现's一致：

  \begin{itemize}
  \tightlist
  \item
    AlN (ACE, 100+ 维) → Theorem 1 can be using ，Theorem 2 non-紧
  \item
    DermaMNIST (SimpleCNN, 128 维but 区分度差) → Theorem 2 紧
  \end{itemize}
\item
  \textbf{两层描述符策略}：SCX 's两层描述符（Layer 1: physical/Domain描述符 →
  Layer 2: Error驱动聚类feature）正Yesat 提升 \(I(\phi; S)\)，will  \(\delta\)
  From大变小。
\end{enumerate}

\begin{center}\rule{0.5\linewidth}{0.5pt}\end{center}

\subsection{附录 A：Symbol表}\label{ux9644ux5f55-aux7b26ux53f7ux8868}

\begin{longtable}[]{@{}lll@{}}
\toprule\noalign{}
Symbol & Meaning & 首次出现 \\
\midrule\noalign{}
\endhead
\bottomrule\noalign{}
\endlastfoot
\(\phi\) & feature映射 & 定义 1 \\
\(\delta\) & 互信息Upper Bound & 定义 1 \\
\(K\) & true state数 & 引理 1 \\
\(\eta\) & 边际Noise率 & Corollaries 2 \\
\(D_m\) & 专家 \(m\) 's损失 & Part 1.4 节 \\
\(C(s)\) & Status一致性 & Part 1.3 节 \\
\(NS(x)\) & Noise分数 & Part 1.3 节 \\
\(\bar{C}\) & 全局平all 一致性 & Corollaries 2.1 \\
\(\varepsilon_\phi\) & 归一化弱度 & 定义 1 \\
\(F1_{\text{base}}\) & 损失基线 F1 & Part 1.4 节 \\
\(F1_{\text{rand}}\) & 随机基线 F1 & Part 1.5 节 \\
\end{longtable}

\subsection{附录
B：关键non-etc.式}\label{ux9644ux5f55-bux5173ux952eux4e0dux7b49ux5f0f}

\textbf{Fano non-etc.式}（引理 1):

\[P(\hat{S} \neq S) \geq \frac{H(S \mid Y) - \log 2}{\log |\mathcal{S}|}\]

\textbf{Pinsker non-etc.式}（定理 2 Proof):

\[TV(P, Q) \leq \sqrt{\frac{KL(P \parallel Q)}{2}}\]

\textbf{Data Processingnon-etc.式}（定理 2 Proof):

\[TV(P_{\text{pred}}, Q_{\text{pred}}) \leq TV(P, Q), \quad \text{where } X \to Y \text{ Yesthe Markov chain}\]

\textbf{AUC 's TV 界}（定理 2 Proof):

\[|AUC_P - AUC_Q| \leq TV(P, Q) \text{ for ArbitraryDistribution } P, Q \text{ 和固定分类器}\]

\begin{center}\rule{0.5\linewidth}{0.5pt}\end{center}

\subsection{References}\label{ux53c2ux8003ux6587ux732e}

\begin{enumerate}
\def\labelenumi{\arabic{enumi}.}
\tightlist
\item
  Fano, R. M. (1961). \emph{Transmission of Information: A Statistical
  Theory of Communications}. MIT Press.
\item
  Cover, T. M. \& Thomas, J. A. (2006). \emph{Elements of Information
  Theory} (2nd ed.). Wiley.
\item
  Pinsker, M. S. (1964). \emph{Information and Information Stability of
  Random Variables and Processes}. Holden-Day.
\item
  Tsybakov, A. B. (2009). \emph{Introduction to Nonparametric
  Estimation}. Springer.
\item
  SCX-Health MedMNIST Experiment报告 v2.
  \texttt{scx-health/results/experiment\_report\_v2.md}.
\item
  SCX AlN v3 Analysis报告.
  \texttt{experiments/mlip\_case/SCX\_AlN\_v3\_Analysis报告.md}.
\item
  SCX 数学Originand Proof建构.
  \texttt{CodexKnowledge/agent\_outputs/05\_数学Originand Proof.md}.
\item
  Theorem 1 (Compression Fidelity).
  \texttt{theory/propositions/04\_compression\_fidelity.md}.
\end{enumerate}

\begin{center}\rule{0.5\linewidth}{0.5pt}\end{center}

\textbf{修正Explanation（2026-06-27）}： 1. \textbf{AUC/PR-AUC \(\eta\)
dependence}：修正了 AUC 和 PR-AUC 界in 遗漏's \(\eta\)
because 子。because this 两个指标涉and conditions概率（分别From \(Z=1\) 和 \(Z=0\)
采样），总变差界is 放大 \(1/\eta + 1/(1-\eta)\) 倍。F1 界non-受影响，because as 
F1 Yes联合Distribution \(P(\hat{z}, Z)\) 'sFunction。 2.
\textbf{Status平衡Assumptions}：as Corollaries 2.2（SCX
退化to 损失基线）添加了Status平衡Assumptions，明确了极non-平衡Statusunder 退化Yes近似's，and 偏差来自
\(\rho(\hat{s})\) 估计误差。 3. \textbf{k-means
Optimal性声明移除}：移除了''k-means 达to  Fano
lower bound'''snot yet 验证声明，改as 存at 性论述加Algorithm近似误差项
\(\varepsilon_{\text{k-means}}\)。

\end{document}
